\section{Winding pressure and long-direction tails}\label{L:part5}

The calibrated winding mass $J_n\asymp n^{-1}$ controls the number of
winding loops per row. We now control how far those loops, and
contractible loops, extend along the cylinder. All polygons in this
section are counted modulo whole-row translations, with the column
labels fixed. Write $h(P)$ for vertical extent in row units.

\begin{proposition}[Exponential cylinder-height tails]
\label{L:prop:cylinder-height-tail}
There are positive constants $c_1,c_2,C$ and $n_0$ such that, for every
regular honeycomb cylinder with $n\ge n_0$ columns and every $u\ge0$,
\begin{equation}\label{L:eq:cylinder-height-tail}
 \sum_{P\ {\rm winding}} X^{|P|}e^{c_1h(P)/n}\le C/n,
 \qquad
 \sum_{\substack{P\ {\rm contractible}\\h(P)>u}}X^{|P|}
       \le C n^C e^{-c_2u/n}.
\end{equation}
\end{proposition}

We obtain the first estimate from the pressure of a gas of disjoint
winding polygons, then join a contractible polygon to a winding one to
obtain the second. Give winding loops fugacity $a=d+d^{-1}$ near
$a=0,d=i$, continuing to kill contractible loops. Let $L_n(a)$ be the
logarithm of the transfer eigenvalue continuing the empty-row-sector
eigenvalue $1$ at $a=0$, evaluated at the physical row parameter.
The analytic part of the proof continues this eigenvalue analytically
along a real interval about zero independent of $n$ and proves
$|L_n(a)|\le C/n$ there. The coefficients of $-L_n(-a)$ are
nonnegative; this converts the real-interval bound into
\[
 |[a^p]L_n(a)|\le C A^p/n \qquad(p\ge1)
\]
with fixed $A$. The factorial in the connected-cluster expansion then
turns these coefficient bounds into the exponential height moment.

\subsection{The continued transfer eigenvalue}

We first identify the continued eigenvalue by a scalar interpolation problem. Keep the pairing states with their annular homotopy, because winding loops now survive with fugacity $a$. The state space is still finite: each cap has only two possible directions around the smaller boundary (displacement less than a full circumference by interlacing of translates). All allowed states can be built from vacant rows at the positive honeycomb point: close an innermost cap first, across its specified interval on the disk side containing no ends, propagating the others as in the finite proof, and repeat. Powers on the nonempty block at zero fugacity are therefore transient by the convergent vacuum bound. Thus the continued eigenvalue and its vector, normalized on the empty component, are analytic locally. Use at first also \(Z\) close to \((1,\ldots,1)\). All rows (spectral parameters \(t\)) commute, by inserting an auxiliary swap and its inverse between the two auxiliary lines at the cut, sliding the swap past each column by braid and canceling around the cylinder; the diagram identities just keep the pairings \emph{with} homotopy when used in the annulus, since the changed patches preserve connections up to isotopy and their internal loops are contractible. Write \(\Lambda(t)\) for the eigenvalue on that same line. It is rational with denominator at most \(D_5 D_6\), \(D_r(t)=\prod_{z\in Z}(z-q^r t)\), using row parity of single corners and the empty row-output evaluation. At \(0,\infty\) a vacant output requires a vacant input and sees the free horizontal straight line or empty, so \(\Lambda(0)=\Lambda(\infty)=1+a\). At zero fugacity \(\Lambda\equiv1\).

For each \(x\in Z\)
\[
 \Lambda(x)\Lambda(q^{-2}x)=\Lambda(q^{-1}x),\qquad
 \Lambda(x)\Lambda(q^{-3}x)=1 .
\]
Indeed in the lower row at \(t=x\) the site is identity, turning its horizontal incoming line up and bottom leg to the right. Thus \(PP^t\) at site \(x\) of the upper row \(t=q^{-2}x\) has a triangle collapsing the two incoming horizontals. Slide this collapse backwards through all other columns (relation \((1\otimes P^t)B_{12}(qs)B_{23}(s/q)=B(s)(P^t\otimes1)\), upper left input numbered first). On reaching the outgoing side of site \(x\), the two triangles there assemble into its cell \(B(q)\) on the fused row. The cap slide at the other specialization works identically, yielding just vertical continuation. This uses the coherent auxiliary half-powers. For generic \(Z\) these equations and endpoint data determine the analytic germ: at linearization at 1 the rational difference with zero value at 0 has basis \(t/(t-q^k z)\), \(k=2,3\). The two rows of linearized constraints at each site have full total rank, as seen for widely separated ordered positive \(z\)'s, where the matrix tends to block triangular with diagonal blocks with entries \(c_k+c_{k+2}-c_{k+1},c_k+c_{k+3}\) in the respective rows, \(c_j=(1-q^j)^{-1}\). Each block determinant is \(((2-\sqrt2)/2)^2\). This proves uniqueness successively on Taylor coefficients.

The following three-term expression is the finite spectral form we will verify; the periodic seam formulation has precedents in \cite[Section~2]{ZhouBatchelor1997}. Use a degree \(n\) polynomial \(U_b(t)\), with roots \(u_j=q^{-1}e^{\theta_j}\). Put \(P(t)=U_b(t)D_4D_5D_6(t)\), using subscripts for shifts of either polynomial at \(q^r t\). The candidate is
\[
 \Lambda= d P_{-2}/P_0+P_{-1}P_2/(P_0P_1)+d^{-1}P_3/P_1.
\]
The two node relations follow immediately using \(P_2=P_3=P_4=0\) at \(t=x\), generically; equivalently the three site-dependent factors apart from \(U_b\) and \(d\) are \(D_2D_3,D_0D_3,D_0D_1\) respectively with common denominator \(D_5D_6\). Endpoints are correct. Cancellation of unwanted root poles (including in \(P_1\)) holds if at \(u_j\)
\[
 D_0 U_{b,-1}U_{b,2}/(D_2 U_{b,-2} U_{b,1})=-d .
\]
We use simple roots with no mutual shifted coincidences.

At \(d=i\), impose that the grade 3 modulo 4 in \(P\) vanish identically. For distinct positive \(Z\) such a solution is given in variable \(s=qt\) by
\[
 \frac{P(t)}{\prod_{x\in Z}(x^4+s^4)}=
 \frac{p_Z^+(-s)}{-s}=1+\sum_x \frac{b_x}{s^2-\sqrt2 x s+x^2},
\]
with \(p_Z^+\) the wave defined earlier (the Gram kernel strictly positive by Fourier). Indeed the quadratic terms on clearing have no forbidden grade, give zeros at \(s=-q^{\pm1}x\), and the interpolation gives the further zero \(s=-x\). Thus the required divisibility holds. Then the candidate is identically 1: after clearing \(P_0P_1\) check at \(0\) and all \(q^j x\). On the orbit starting at \(x\), \(P_2=P_3=P_4=0,\ P_0=iP_6,\ P_7=i(P_5-P_1)\) by grade elimination, giving equality throughout. Confluence is allowed since the wave exists there by the confluent positive Gram; the leading coefficient in \(s\) of \(P\) stays 1.

\begingroup\emergencystretch=1em
The remaining \(n\) zeros in \(s\) are positive simple at distinct positive sites. Here are details by variation. This is a total-positivity and variation-diminishing argument in the framework of Schoenberg \cite[Theorem~1 and p.~17]{Schoenberg1947}; we prove the strictness and decay needed for the present kernel. In log coordinates, convolution against the kernel with Fourier multiplier \(\sinh(\pi\nu/4)/\sinh(3\pi\nu/4)\) takes \(-p_Z^+(-e^y)\) to \(p_Z^+(e^y)\). Indeed it acts identically on the exponential (contour shift in Fourier inversion gives exponential decay with rate \(>1\), so use the multiplier at imaginary frequency) and takes \(-M(-e^y,x)\) to \(M(e^y,x)\), by the Fourier transform of \(\sin b/(\cosh y-\cos b)\), which is \(2\pi\sinh((\pi-b)\nu)/\sinh\pi\nu\) by residues. This kernel is strictly totally positive. To recall a proof, by the sine product (equivalently integrate the integer-pole cotangent partial fractions for its log derivative) it is proportional to the convolution of normalized double exponentials of rates \(4k/3, k\ge1, 3\nmid k\) (the independent centered series converges by variances). Each factor has nonnegative minors on ordered tuples: decompose into two one-sided exponential kernels, whose minors reduce to step matrices, and integrate determinants. Infinite tails inherit nonnegativity by convergent convolution, for example by dominated Fourier inversion retaining several factors. Strict positivity at order \(l\) follows by convolving sufficiently many initial factors: the double exponential has strict determinant between aligned ordered tuples differing by motion of one point without passing neighbors (in composing step matrices there is then an open set of strict interlacings). Any two tuples can be joined by at most \(2l\) such steps via intermediate points far off, with strictness on open sets. The final tail determinant into a given tuple is positive somewhere (linear independence of translates by nonzero Fourier transform). The determinant integrations here are just antisymmetrization on products of ordered lines. This proves the assertion. If the input of our convolution had fewer than \(n\) sign changes, a linear combination of translates at some of the \(n\) output zeros \(\log x\) could match its signs: use the determinant of that many translates (one more than the number of changes) evaluated at the change points and the running point. Strictness and ordering give the signs, contradicting zero integral (the kernel decays faster than required for the exponential). Thus \(-p_Z^+(-s)\) has at least \(n\) sign changes, proving the root claim.
\par\endgroup

\subsection{Counting phases and sparse end roots}

We detail the analytic control of the root equations. The target is an inverse bound uniform in the circumference; in particular the first and last roots cannot be treated by a bulk Riemann-sum approximation. For positive sites put
\[
 A_b(x)=-\arg(1-e^{ib}e^x)/\pi,
 \quad A=A_{\pi/4},\quad K=A-A_{\pi/2},\quad k_b=K'.
\]
Arguments are continued from 0 at the left end. Thus
\(2\pi k_b(x)=t_0/(\cosh x-t_0)-1/\cosh x,\ t_0=2^{-1/2}\). At \(d=i\) order the real roots increasingly; their pole equations are
\[
 F(\theta_j)=j-1+3/4,\qquad
 F(x)=\sum_{z\in Z}A(x-\log z)+\sum_j K(x-\theta_j),\quad 1\le j\le n .
\]
Indeed the fractions give these phases modulo 1 and the labels are constant by continuity on ordered distinct sites. At widely separated sites the wave formula has one root \(e^{\theta}\sim(1+\sqrt2)z\) at each scale: in the displayed quadratic fractions the zero equations at \(s=-z\) give a triangular limiting system on \(b_x/x^2\) with diagonal \(1/(2+\sqrt2)\) and entries 1 for larger columns, hence nonzero solutions there and roots on bounded site ratios by substitution. Then \(A(\log(1+\sqrt2))+K(0)=3/4\), previous sites contribute 1 each. These equations persist at homogeneous confluence (bounded coefficients with fixed leading term by the Gram; no roots at 0 in the limit since the constant fraction there is the homogeneous wave limit \(p_Z^+(-s)/(-s)|_{s=0}=(-1)^n Y_n\ne0\) proved above). Different labels imply simplicity there too. For nearby \(d\) the targets are shifted uniformly and analytically by the phase increment divided by \(2\pi\).

Hereafter at homogeneous confluence the continuum comparison has cumulative \(H(-\infty)=0\),
\[
 h_b(x)=H'(x)=\frac{4}{\sqrt3\pi}\frac{\sinh(8x/3)}{\sinh(4x)},\qquad H=A+k_b*H
\]
(by the same Fourier formula, multiplier for \(h_b\) is \((2\cosh(\pi\nu/4)-1)^{-1}\)). Define
\[
 E(x)=\#\{j:\theta_j\le x\}-n H(x)+1/3,\qquad B_j=n h_b(\theta_j).
\]
At jumps \(E\) takes one-sided values \(\pm1/2+(k_b*E)(\theta_j)\), since \(F=nH+k_b*(E-1/3)\) and \(\int k_b=1/4\). Monotonicity between then gives \(\|E\|_\infty\le (1/2)/(1-\|k_b\|_1)<.711\).

The linearization of the root equations is the map
\[
 (\mathcal L_n\eta)_j
   =F'(\theta_j)\eta_j-\sum_l k_b(\theta_j-\theta_l)\eta_l.
\]
Its natural displacement norm is
$\|\eta\|_B=\max_j B_j|\eta_j|$: a root in a sparse tail is allowed
a larger displacement than one in the bulk.

\begin{lemma}[Uniform continuation of the roots]
\label{L:lem:pressure-root-continuation}
For all sufficiently large $n$, the homogeneous roots at $d=i$ satisfy
$\min_jB_j\ge b_*>0$ and
\[
 \|\mathcal L_n^{-1}v\|_B\le C\|v\|_\infty.
\]
There is a fixed disk about $a=0$ on which they continue analytically
as simple roots, with
$|\theta_j(a)-\theta_j(0)|\le C|a|/B_j$.
\end{lemma}
\begin{proof}
The only difficulty in a bound uniform in $n$ is the first few roots
at either end. We first record bounds on the kernel, then locate these
roots in cumulative coordinates, and finally verify the inverse bound
by diagonal dominance.

\paragraph{Kernel bounds and root locations.}
Useful elementary bounds are
\[
 \int(k_b)_+<.2725,\quad\|k_b\|_1<.296\ \text{(in fact }<.295),\quad
 \operatorname{Var} k_b<.475,\quad -.0058<k_b<.226.
\]
Here \(\operatorname{Var}\) denotes total variation. \(k_b\) is positive for \(|x|<d_0=\operatorname{arcosh}(1+\sqrt2)\), and decreases on the positive half until a unique minimum (differentiate in \(\cosh x\)). The first bound and \(\|k_b\|_1<.295\) follow from \(K(d_0)<.26125\) by the displayed argument formula; more precisely \(K(d_0)<.2612\), which gives the stated bound .711 above. The negative part satisfies \((k_b)_-\le .094 e^{-|x|}\) directly. Also with \(c=3/4\), \(e^{d/c}(k_b)_+(d)\) is unimodal with maximum \(<.279\) and
\[
 A_p'=\int_0^\infty(k_b)_+(c\log v)dv<.58 .
\]
For verification, on \(0<d<d_0\)
\(-k_b'/k_b=(1-t_0)\sinh d/(t_0-(1-t_0)\cosh d)+\sinh d/(\cosh d-t_0)+\tanh d\); it crosses \(1/c\) once, between .27 and .29 (increasing up through \(\cosh d=1/t_0\), then staying greater), with \(k_b(.27)e^{.29/c}<.279\). It also exceeds \(c^{-1}\tanh(d/c)\), e.g. bound the middle and last terms alone separately for \(d\le .5\) and beyond. Thus \(k_b\cosh(d/c)\) decreases, and its integral up to \(d_0\), times \(2/c\), gives the last inequality (upper sum with step .4 and left values bounded by .226,.177,.098,.037 respectively). Here are additional details for checking these numerical inequalities if desired. With \(r=e^{d_0}=1+\sqrt2+\sqrt{(1+\sqrt2)^2-1}\),
\[
 K(d_0)=\tfrac14+\pi^{-1}\arctan\frac{(2t_0-1)r-1}{1-r+2t_0 r^2}<.2612
\]
using \(4.61<r<4.62\); the minimum of \(2\pi k_b\) is \(-(1-\sqrt{t_0})^2/t_0\). In the crossing test the sum at .27 is \(<1.31\), at .29 is \(>1.36\); on \([.5,d_0]\) the middle term alone is \(>1\) (check endpoints), last term \(>.4\), while on \((0,.5]\) the middle term exceeds \(d/c^2\). The displayed point-value checks follow by substitution (e.g. \(3.1415<\pi<3.1416,\ .70710<t_0<.70711\); for hyperbolics at \(0<x<2\) sum through degree 12 for \(\cosh\), 13 for \(\sinh\), with errors \(<.000001\) by the Taylor series). These relaxed constants will suffice.

From the jump bounds, right-tail cumulative coordinates \(u_l=n(1-H(\theta_{n-l}))\) are \(l+1/6\) up to error \(<.211\). Moreover \(u_0>.097\): in the convolution at the last root use \(E\ge1/3\) on the right half, giving
\(u_0\ge1/6-(.711/2)\|k_b\|_1+(\int(k_b)_+)/6-(.711/2)\int(k_b)_-\).
Left coordinates \(u_l=nH(\theta_{1+l})\) are similarly \(l+5/6\) up to .211. Thus \(B_j\) stays positive; at finite depth from either end it is asymptotically \(u_j/c\) using that end's coordinates (reindex from 0). Distances on one end tend for bounded indices to \(c|\log(u_l/u_j)|\), up to \(o(1)\). We use subsequences to fix all such limits.

\paragraph{Diagonal dominance in the bulk and at the two ends.}
After rescaling columns of the linearization by \(1/B_l\) it suffices to show strict row diagonal dominance with uniform margin. We can subtract the self term \(k_b(0)\), or equivalently include it in the row sum to be compared with \(F'/B_j\). Always \(F'(\theta_j)\ge B_j-.711\operatorname{Var} k_b\). If depth from both ends tends to infinity, the kernel absolute row sum tends to \(\|k_b\|_1\), by bounded discrepancy and Riemann approximation near that root. Here and below far-position tails converge uniformly for exponentially decaying entries: each unit \(x\)-bin has total sampling weight \(\sum 1/B_l\le C\), by bounded discrepancy, the density changing by bounded factors, and the lower bound on occupied \(B_l\). At fixed depth the row sum thus takes the limiting form
\(\sum_l c |k_b(c\log(u_l/u))|/u_l\), \(u\) the root's coordinate.

For a limiting tail root with \(u\ge .955\), the positive part of this sum is bounded by
\(.2725+2(.711)c(.279)/u\): integrate against unit density in the cumulative coordinate and bound the quadrature error by .711 times the variation (one may center the discrepancy between samples and volume by a constant, by the indicated ranges). The upper-ratio negative tail \(u_l/u>e^{d_0/c}\) costs at most
\(.0113+2(.711)c(.0058)(.131)/u\) similarly (positive variation of the single negative hump times the decreasing factor is bounded by its rise times the largest factor, \(e^{-d_0/c}<.131\)). On \(u_l<.131u\) bound the first sample separately by \(c(.0058)/.097\); the others cost together at most
\[
 .094\,(.955)^{c-1}(.131+.045/u)^c
\]
using \(u_l\ge l-.045\), summing \(c\,u_l^{c-1}u^{-c}\). The total is strictly less than \(1-.711(.475)c/u\). 

For the first sample at the left end (if \(u<.955\); in any case \(.622<u<1.045\)), the positive sum is at most \(c(.226)/u+c(.226)/1.622+.2725/2\), by the decreasing integrand on \(u_l\ge l+.622, l\ge1\) and the half integral. Only the upper negative tail is needed, as just bounded. This still gives strict dominance.

Finally consider the first right sample \(.097<u<.378\). The off-self sum is bounded by \(c(.09)/.955+.2725/2\) plus the same upper negative tail (the positive term uses \(k_b(c\log(.955/.378))<.09\)). Here improve the diagonal using
\(F'=B_j+\int k_b(\theta_j-x)\,d(E(x)-1/3)\).
After deleting self, the integral correction has limit the quadrature difference for \(k_b(c\log(v/u))\) on \(v>0\), samples \(u_l,l\ge1\) minus unit density. This difference converges despite the individual power tails, by integration by parts with bounded discrepancy (also controlling the far end before taking limits, since \(k_b,k_b'\) decay exponentially in position distance). The omitted volume up to \(v=2/3\) costs at most \(.58u\); on \([2/3,\infty)\) discrepancy counting from that point is bounded by .711. Variation there is at most \((.279/(2/3))u+.0116\) by monotonicity up to the minimum. Thus the limiting physical diagonal after self deletion is at least
\[
 u/c-.58u-.711((.279/(2/3))u+.0116).
\]
After division by \(u/c\) this too strictly exceeds the stated off-self bound. This proves the uniform inverse claim.

\paragraph{Analytic continuation.}
These estimates give an analytic root germ on a disk of fixed radius in the uniform target shift. Indeed perturb by \(\eta_j\); the left sides become
\[
 F(\theta_j+\eta_j)+
 \sum_l [K(\theta_j-\theta_l+\eta_j-\eta_l)-K(\theta_j-\theta_l+\eta_j)].
\]
The analytic Taylor remainders are uniformly quadratic in the weighted norm on a small ball: near each root positive-order derivatives of \(F=nH+k_b*(E-1/3)\) are controlled on fixed small strips by \(C(B_j+1)\), and the difference-sum uses exponentially decaying kernels with at least one factor \(1/B_l\), summable by the bin estimate. Applying the inverse and contraction therefore gives \(|\eta_j|\le C|a|/B_j\). Simplicity under sufficiently small shifts follows too (spacing locally bounded below in inverse density units by the level equations). This proves the lemma.
\end{proof}

For fixed large \(n\) this solution germ also exists at sites sufficiently near homogeneity; at \(d=i\) it coincides with the grade-elimination solution above by continuity and local uniqueness. Thus by the scalar uniqueness for generic sites and then continuity, the continued eigenvalue is given by our candidate. Root poles cancel as polynomial factors before evaluation (also at an accidental specialization of a node), so evaluations away from \(D_5D_6=0\) are analytic on the small homogeneous root disk.

\subsection{The pressure error at homogeneous sites}

The uniform inverse has constructed the eigenvalue on a fixed parameter neighborhood. We now need its logarithm to be of order the reciprocal circumference. We sharpen the estimate at homogeneous sites to order \(1/n\). We need only real phases \(d=e^{i\phi}\) close to \(i\), so the deformed \(\theta_j\) are real by the construction. Now let \(N(x)=\#\{\theta_j\le x\}\) for these roots, and \(F\) also denote the deformed counting phase. Put \(\ell=\phi/(2\pi)-1/2\), so \(F(\theta_j)=j+\ell\). The shift estimate and exponential density bounds imply \(N-nH\) remains uniformly bounded. With \(C_\phi=-(4/3)\phi/(2\pi)\),
\[
 F=nH+k_b*(N-nH),\qquad (1-k_b*)(N-nH-C_\phi)=N-F+\ell+1/2=:r .
\]
We will estimate the change of logarithmic root products by integrating
against this discrepancy. The needed cancellation is stronger than its
uniform boundedness:
\begin{lemma}[Quadrature against the root discrepancy]
\label{L:lem:pressure-quadrature}
For the real root deformations just constructed, let $p$ be a smooth
test whose value and first two derivatives are bounded by
$D\exp(-4|x-b|/3)$. Uniformly over the fixed real phase interval,
\[
 \left|\int p r\right|\le C D e^{4|b|/3}/n .
\]
\end{lemma}
\begin{proof}
In the interval \(|x|\le(3/4)\log n-C_0\), with sufficiently large fixed \(C_0\), we have \(F'\asymp n h_b\), and derivatives up to the third bounded by \(C n h_b\). There \(r=1/2-\{F-\ell\}\) a.e., since both neighboring level roots lie in the increasing range by bounded discrepancy (take it increasing also within unit distance of the interval). Integrate by parts twice using bounded periodic zero-mean primitives of the mean-zero sawtooth, dividing by \(F'\) each time. Boundary terms cost \(C D e^{4|b|/3}/n\); the remaining integral costs at most \(C D e^{4|b|/3}\int_{|x|\le(3/4)\log n} n^{-2} e^{4|x|/3}dx\). Outside that interval of integration by parts the estimate is direct by decay.
\end{proof}

\paragraph{Dominant spectral terms.}
Use spectral parameter \(qt=e^z\), and call the three summands of the candidate \(T_1,T_2,T_3\) respectively. Their consecutive ratios are \(\exp(2\pi i F(z))/d\) and \(\exp(2\pi i F(z+i\pi/4))/d\) by the phase formulas (here arguments are in log-coordinates). On the line \(z=b+i\pi/8\), \(b\) real, terms 2 and 3 relative to 1 are bounded by
\[
 C\exp(-c' n e^{-4|b|/3})
\]
for fixed \(c'>0\). On \(z=b-i\pi/8\) terms 1 and 3 relative to 2 satisfy the same bound. To verify this, take log moduli and replace root sums by \(n h_b\) integrals (here \(h_b\) still denotes the continuum kernel above, independent of the horizontal shift \(b\)). Errors are \(O(1)\) by discrepancy, kernels differentiated along the horizontal having uniform exponential decay about \(b\). At argument \(w=b\pm i\pi/8\), the continuum phase per column is \(H(w)=A(w)+\int K(w-y)h_b(y)dy\); its imaginary part has the desired signed lower bound. Indeed it vanishes at spatial infinities, and the horizontal derivative has imaginary part that of \(h_b(b\pm i\pi/8)\), equal to \(\mp(2/(\sqrt3\pi))\sinh(8b/3)/\cosh(4b)\). For \(T_3/T_1\) on the upper line one adds the phase at \(w+i\pi/4\). The summed continuum phase still has imaginary part tending to zero at the ends, but now derivative \(h_b(w+i\pi/4)\). Indeed continue the differentiated identity by lifting the convolution contour upward to keep the first pole at \(y=(w+i\pi/4)-i\pi/4=w\) below it (lift by less than \(\pi/4\); isolated meromorphic poles are harmless). Returning to the real contour adds the residue correction \(-h_b(w)\) to this continued formula since the residue in \(y\) of \(K'(w+i\pi/4-y)\) is \(-1/(2\pi i)\). The resulting imaginary derivative of the sum is \(-(4/(\sqrt3\pi))\sinh(8b/3)/\cosh(4b)\), again sufficient.

The dominant summand (1 above, 2 below) is itself bounded uniformly. Indeed the same integral substitution with \(O(1)\) error applies to its log modulus; the root factors per root are smooth log ratios of \(e^y-e^{b+i\alpha}\), with angles a fixed distance from zero modulo \(2\pi\). The resulting thermodynamic term per site must vanish identically: at every fixed \(b\) use \(d=i\), where the eigenvalue sum equals 1 and the relative bounds just proved tend exponentially to zero.

\paragraph{The pressure change.}
Compare that dominant term to itself at \(d=i\) as the roots deform, writing \(g(y)\) for the complex log of its root factor tending to zero at \(+\infty\). It is a sum of one or two log differences as above, with angle differences (angles taken in \((0,2\pi)\)) as follows:
\begin{itemize}
\item Upper: numerator from \(z-i\pi/2\), denominator from \(z\), giving difference \(3\pi/2\).
\item Lower: numerators from \(z-i\pi/4,z+i\pi/2\), respective denominators from \(z+i\pi/4,z\), differences \(3\pi/2,-3\pi/2\).
\end{itemize}

The dressed derivative \(p=(1-k_b*)^{-1}g'\) then satisfies the quadrature hypothesis uniformly. Indeed in a Fourier convention with \(e^{i\nu(y-b)}\), a single difference's transformed derivative is \(2\pi i(e^{-\nu\alpha}-e^{-\nu\beta})/(1-e^{-2\pi\nu})\) by shifting up a period. Resolving the convolution replaces the denominator's \(\sinh(\pi\nu)\) by \(\sinh(3\pi\nu/4)(2\cosh(\pi\nu/4)-1)\). The double zeros at \(\pm4i/3\) then yield only simple poles, by the angle differences. Exponential horizontal decay and contour shift beyond these first poles prove the claim, including derivatives. Thus against \(N-nH-C_\phi\) the integral of \(g'\) equals that of \(p\) against \(r\) by even convolution. The change of log root products is minus the integral of \(g'\) against the change of \(N\). Its constant contribution \(g(-\infty)(C_\phi-C_{\pi/2})\) exactly cancels the explicit phase change in \(T_j\) (\(g(-\infty)\) is \(3\pi i/2\) above, zero below). The remaining log change costs \(C e^{4|b|/3}/n\).

Altogether, using the bounded dominant term directly if this last bound is large, we have \(|\Lambda(t)-1|\le C |\cosh(4z/3)|/n\) on the two lines. The same holds between by maximum modulus on division: the rational function in \(e^z\) has bounded endpoints and no poles there. Consequently \(L_n(a)\) continues analytically along a fixed short real interval about zero, with modulus \(O(1/n)\) there (the simple real root construction extends locally analytically even when an evaluation aligns with a removable pole).

\subsection{Connected clusters and exponential height tails}

We explain the geometric consequence. The pressure coefficients count connected families of winding polygons; their common sign will turn the analytic bound into positive height moments. On a finite long strip of this cylinder with vacant boundary input and output the partition function is a gas of pairwise disjoint winding polygons. Its log per row coefficients tend to \([a^p]L_n\) by simplicity/dominance locally at zero fugacity. The connected-cluster formula at order \(p\) sums products of the \(X^{\rm length}\) weights divided by \(p!\), with multiplier the alternating subgraph sum over connected spanning subgraphs of the overlap graph (also conflicts for repeated polygons). This follows by expanding disjointness pair factors and taking the formal log. This is the Mayer connected-graph expansion; see \cite[Propositions~2.1, 2.6 and~2.8]{ScottSokal2005} for the expansion and its sign in a repulsive gas. Nonzero multipliers have common sign \((-1)^{p-1}\) and modulus at least 1, for example by deletion-contraction for the linear coefficient of the graph coloring polynomial. Thus we may lower bound absolute coefficients even in infinite length by restricting to stars where all other polygons meet the first (first counted per unit row, by the finite-strip limit).

The passage from a real interval to a coefficient bound is the positive-coefficient principle used in \cite[Section~2.3]{ScottSokal2005}. This sign property also completes the coefficient estimate: the power series for \(-L_n(-a)\) has nonnegative coefficients. Its radius must reach across our short interval by analyticity along that interval (otherwise at the first positive convergence boundary all derivative series still sum to the derivatives by monotonicity, and Taylor expansion there gives positive convergence farther out). Evaluating inside then bounds each fixed coefficient by \(O_p(1/n)\).

We next turn a coefficient into a height moment. Given a winding
polygon $P$, the weighted choices of winding loops meeting it satisfy
\begin{equation}\label{L:eq:winding-overlap}
 \sum_{Q:\,Q\cap P\ne\varnothing}X^{|Q|}
       \ge (h(P)-C)_+J_n,
\end{equation}
where $Q$ here runs over actual row positions. To prove this, choose
vertices $p_+,p_-$ on the same bipartite sublattice, within bounded
distance of the highest and lowest levels of $P$. Their row-height
difference is an integer $D\ge h(P)-C$.

For a second winding polygon $Q$, write $Q_{b,r}$ for its translate by
column phase $b\in\mathbb Z/n\mathbb Z$ and row shift $r\in\mathbb Z$.
Define the signed mean height $m(Q)$ by integrating the difference
between the indicator of its lower domain and that of the lower
half-cylinder, divided by one period's area per row. Column translation
preserves $m(Q)$, and $m(Q_{b,r})=m(Q)+r$. For disjoint essential curves,
mean-height order agrees with their nesting order, since the area
between them is positive.

Put $t=m(P)-m(Q)$. If $r>t$ and $Q_{b,r}$ is disjoint from $P$, then
$p_+$ belongs to its lower domain. If $r-D<t$ and $Q_{b,r-D}$ is
disjoint from $P$, then $p_-$ belongs to its upper domain. For each
integer $t<r<t+D$, average these two necessary conditions over column
phases. The relative query points have the same distribution: their
row coordinates agree after subtracting $D$, and their common bipartite
type makes the remaining displacement an integer column shift. The
upper and lower open domains of $Q$ are disjoint. Hence the two
probabilities of meeting $P$ sum to at least one.

There are $D-O(1)$ such integers $r$. Their upper shifts $r$ and lower
shifts $r-D$ form disjoint intervals, so summing over all shifts gives
at least $D-O(1)$ intersecting translates in the phase average. Finally
sum $Q$ modulo row translation with weight $X^{|Q|}$. Column translation
permutes this weighted family, so its phase average is the original
sum, with no additional factor of $n$. This proves
\eqref{L:eq:winding-overlap}.

Let $c_{n,p}=|[a^p]L_n(a)|$. Evaluation of the nonnegative series at a
fixed positive argument $a_*$ inside the interval gives
$c_{n,p}\le C a_*^{-p}/n$ for every $p\ge1$. Restrict the ordered
connected-cluster sum to tuples whose first polygon meets every other
one. Their overlap graphs are connected, and the multiplier has modulus
at least one. For $m\ge0$, the factor $1/(m+1)!$ in that expansion gives
\[
 (m+1)!c_{n,m+1}
 \ge\sum_PX^{|P|}
       \left(\sum_{Q:\,Q\cap P\ne\varnothing}X^{|Q|}\right)^m
 \ge J_n^m\sum_PX^{|P|}(h(P)-C)_+^m.
\]
Since $nJ_n$ is bounded below, choose $\varepsilon>0$ so that
$\varepsilon/(a_*nJ_n)$ is uniformly less than one. Divide the last
inequality by $J_n^m$, multiply by $(\varepsilon/n)^m/m!$, and sum.
The exponential moment is bounded by
\[
 \frac{C}{a_*n}\sum_{m\ge0}(m+1)
      \left(\frac{\varepsilon}{a_*nJ_n}\right)^m=O(n^{-1}).
\]
The harmless factor $e^{\varepsilon C/n}$ restores $h(P)$ in place of
$(h(P)-C)_+$. This proves the winding estimate in
\eqref{L:eq:cylinder-height-tail}. It also gives suppression beyond
$n^{1+\epsilon}$ faster than every fixed power, even with fixed length
powers: deterministically $|P|=O(n(h(P)+1))$.

The contractible height tail follows by a local joining operation. Join any such polygon to an arbitrary winding one by shifting the latter to put a bottom extremity just above a top extremity of the contractible one. More explicitly, using horizontal triangle-side levels, a highest triangle-center vertex must be down-pointing, using both lower-slant connections. Take the hexagonal face for which it is the lower left vertex just below mid-level. Intersection with the lower boundary of this face is then precisely a single arc of one or two edges (if the other extremity-level vertex on it is used, it too must use the shared lower neighbor). Align the winding polygon by lattice translations to have a lowest vertex (up-pointing) on the corresponding upper left of that face, giving the symmetric upper arc. They are otherwise disjoint, so flipping the two arcs to the complementary face edges gives a single winding loop with bounded length change and retaining the height up to bounded loss. Reversing from an output costs at most order its length in face choices and an additional factor \(O(n)\) to restore the original horizontal phase. Thus multiplying the contractible tail mass by \(J_n\) bounds it by a polynomially weighted winding tail, sufficient again with length powers allowed.


To obtain the second estimate of
\eqref{L:eq:cylinder-height-tail}, apply the winding exponential moment
with a smaller exponent, absorbing the polynomial recovery factors and
the deterministic bound $|P|=O(n(h(P)+1))$. Division by
$J_n\asymp n^{-1}$ introduces only another polynomial in $n$.
The bounded change of height in the join changes the constant.
This completes the proof of Proposition~\ref{L:prop:cylinder-height-tail}.
